% Proposed full-depth accessible insertion; not a summary of the technical text.
% Inherit the selected master's Noto Sans/Fira Math setup.
\section{More data must survive a harder test}

Give a model another thousand examples of the same reach and it may get better at that reach. Give it examples from another person, another task or another kind of body and you change what it has to learn. That is the attraction of a Large Movement Model: the relationships learned from a varied collection might help it cope with a movement it has never encountered. It is also a question we can get wrong. A broader collection could help, add nothing or teach a correspondence that makes a particular prediction worse.

The proposal therefore makes two bets that need separate tests. One is about amount: does more qualified training data improve the chosen prediction? The other is about diversity: does a wider range of qualified examples help when their total amount is held comparable? Neither bet has been demonstrated here. No body-wide world model has been trained or validated by this proposal.

\subsection{Four source groups, distinct observations}

The proposed training collection has at least four groups. One supplies neural observations. A second supplies muscle and other myographic observations: records of muscle activity whose actual measurement type must be stated. A third supplies measured movement and displacement. A fourth supplies examples generated by physics-based motion or musculoskeletal simulation. Each would teach the model something different. Listing them does not mean that four complete collections have been acquired, or that nothing else could be useful.

The intended neural range includes suitable brain, spinal, peripheral and relevant visceral recordings. A recording's site, sensor and role in the task determine what it can tell us. The current focus is motor--sensory. The programme uses Myoception for sensing muscle-driven movement and the feedback that comes with it, including body-position and movement feedback, touch and other task-relevant sensations. This proposed programme use of the name establishes neither a new external scientific standard nor a simultaneous record of the whole body. The intended range remains broader than the available paired examples.

Touch, body-position and movement feedback (proprioception), contact, force and context cannot be treated as interchangeable labels. A known tactile stimulus says what the experiment delivered. A contact sensor reports the event it was built to detect. A calibrated force instrument measures a specified push or pull. Movement measurements can help estimate where a limb is without directly recording the sensory nerve signals or how its position feels to the person. EMG records muscle electrical activity; it does not directly record peripheral nerve or sensory activity. A stimulus, a report of sensation, a recorded neural feature and a model's estimate therefore keep their different meanings even when they concern the same action.

The simulation group is a training source in its own right. A qualified virtual body could supply examples of motion, changes between states, calculated muscle--tendon and contact forces, or movement patterns used as a starting point for learning. Each example would retain the simulator and body-model version, parameter settings, numerical solver, controller and its original sources. A controller that requests a joint position differs from one that drives muscles. Calling both muscle control would conceal information about how the movement was produced.

If a measured motion is transferred to another body model, keep the original recording and the transformation with it. If the computer then calculates internal states or forces, those values remain simulated. A calculated contact force does not become an independently measured one because a measured trajectory helped produce it.

The generator first needs calibration against suitable measured support, with uncertainty and data roles kept explicit. Its adjustable settings would use training support, leaving independent tests untouched. A proposed comparison would keep a fixed amount of real data in both arms and add simulation examples in one arm. It would match the permitted information, computing work and original decision rules, then check the contribution against held-out real observations. Record how often each arm sees real and simulated examples and how many updates it makes. Keeping the same real recordings does not ensure that the models see them equally often. Count simulated examples and their variety separately from real acquisitions. Real test outcomes must never be used to tune the generator, calibrate it, select rewards or choose the model that will face those same tests.

Simulation could teach a useful relationship, or pass on its own wrong contact rules, dynamics or body correspondence. Those failures must remain visible. Generated examples cannot independently verify the biological assumptions that generated them. The methods, implementation and use conditions of any proposed generator still need checking; no frontier implementation or dataset is selected here.

These four source groups do not replace the programme's five component roles. Brain signals, muscle signals, measured movement, a virtual body and a learning model describe what each part does in the system. The source groups describe what can supply training examples. The virtual body can supply examples and constraints, while the learning model can use several groups. Other useful information can be added if its meaning and measurements can be checked.

\subsection{Count what the model actually receives}

A recording can be divided into many short windows, and the same window can be shown to a model many times. Those operations increase the number of training presentations. They do not create another person or another independent recording. For every input type, the study would report the original acquisition groups, the people or animals, recording time, trials, available target values, distinct windows and repeated presentations. An acquisition group keeps together material that came from the same underlying recording, including copies, overlapping extracts and differently named versions.

Diversity needs a description too. Different people, recording days, movements, body regions, sensors, apparatus and environments expose different changes. A dataset containing one new sensor is not equivalent to a dataset containing twenty new people. One more dataset is therefore an inadequate measure of breadth. The proposed study would identify which aspect changes and which aspects remain comparable, including the quality and missingness of the measurements.

To test amount, we would increase eligible training material from a declared source distribution. To test diversity, we would hold a stated amount comparable while changing selected aspects of the collection. A design that crosses both changes could ask whether their effects depend on each other. The model design, permitted inputs, prediction target, preprocessing and selection rules would remain fixed. The experiment would report how often examples were reused, the number of updates made to the model and the computing work actually consumed. Equal numbers of updates need not cost the same or expose the model to the same information.

Some comparisons will resist a clean match. Adding a sensor type may also change the task, participant population or recording quality. In that case we would report the combined change and its uncertainty. Calling it diversity alone would assign the result a cause the comparison has not separated.

These are proposals for a later, separately specified study. Before running it, the exact recordings, training budgets, random starting conditions, controls and pass/fail rules would have to be fixed. The proposal does not change the already adopted continuous WAY-EEG-GAL development comparison, its participant roles or its prescribed updates. It does not open reserved LINK recordings, ST or other protected tests. Calibration and numerical scaling would use training support only. A recording already seen by the model cannot be made unseen by giving its windows new filenames.

\subsection{Ask where the improvement holds}

A new recording day from a familiar person tests something different from a new person. A new person tests something different from a new movement, a new apparatus or another species. Generalization means carrying a learned relationship into a stated new setting, so each of those changes needs its own result. The study would keep the relevant original recordings together before making windows or choosing settings, and would audit any earlier pretraining exposure.

Each proposed transfer would face a model trained on target data alone, the applicable deliberately incorrect correspondence and appropriate simpler predictors. The original tests for false associations, shuffled pairings, shared cues, recording artefacts and access to future information would remain. Comparing a source-trained model only with an untrained model would not show whether the source data adds useful information beyond a simpler alternative. Nor could a better average erase harm to a recording or output protected by an original no-harm rule. Unhelpful transfer, plateaus and outright failures belong in the result.

\subsection{What a power law would mean}

One possible pattern is that error falls as training amount rises, with each further increase buying a smaller improvement. A proposed mathematical description is
\[
E(N)=E_{\infty}+A\left(N/N_0\right)^{-\alpha}.
\]
The letters name quantities that the study would define before fitting the curve. $N$ is the chosen training amount and must be greater than zero. $N_0$ is a fixed positive reference amount in the same units, so dividing one by the other leaves a pure number. $E(N)$ is a nonnegative error on the specified held-out population. For this proposed decreasing family, the floor $E_{\infty}$ must be zero or positive, the amount-dependent size $A$ must be zero or positive, and the exponent $\alpha$ must be positive. The floor and $A$ have the error's units; the exponent has none. If $A$ is positive, this curve falls as training amount increases. If $A$ is zero, it stays flat however large the positive exponent. The exponent alone therefore cannot establish the direction. If the error measure has an upper limit, the predicted values must stay inside it over the declared fitting and prediction range as well. A fitted floor would describe this candidate curve; it would not establish an irreducible limit of biology.

There is no estimated exponent or demonstrated law here. Choosing a decreasing candidate does not require the measurements to improve. A flat line, a plateau, rising error, a reversal or a different pattern for each domain may describe them better. Harm from added data must be reported, even when the proposed curve cannot represent it. The study would choose its fitting range, alternative curve shapes, uncertainty method and checks on withheld data-amount levels in advance. It would ask whether the fitted relationship predicts levels left out of fitting, rather than merely drawing a pleasing line through the points it has already seen. Overlapping windows from one recording do not provide the independent repetitions needed to make uncertainty vanish.

The measurement being fitted also stays fixed. An error in rescaled muscle-signal amplitude, an error in camera pixels and an error in physical joint angle describe different things. Combining them into one percentage called biological accuracy would need a defensible, independently checked scoring rule. The proposal instead retains the original units and reports each domain, the number of observations, absent values and uncertainty, alongside any justified common summary. More training data, a larger computational model and averaging over more outputs remain different changes.

\subsection{Two useful sources do not make one shared event}

FALCON M1-A illustrates a possible neural--muscle link. The fixed release describes simultaneous motor-cortex activity and processed electrical activity recorded with electrodes within the muscles of one rhesus macaque. That acquisition is called intramuscular EMG; it differs from a surface recording taken over the skin. Processing the signal is a separate step from where it was recorded. Its three split labels do not mean three animals. The reviewed metadata does not establish a released synchronized motion-capture trajectory. Some processing of the muscle targets uses information on both sides of a time point, including later information. A model scored against those targets would therefore need a separate test before being described as purely causal in real time. Its small minival subset is drawn from calibration data to check submission format. Because those examples are already part of calibration, they are not an independent validation population. \href{https://doi.org/10.48324/dandi.000941/0.241029.1405}{FALCON M1-A release}; \href{https://papers.nips.cc/paper_files/paper/2024/file/8c2e6bb15be1894b8fb4e0f9bcad1739-Paper-Datasets_and_Benchmarks_Track.pdf}{FALCON methods}.

AJILE12 supplies a different candidate link. It connects recordings inside the skull with upper-body positions estimated from video in twelve clinical participants. The positions are two-dimensional image coordinates, measured in pixels, rather than instrumented three-dimensional motion capture. The reviewed release does not establish paired muscle electrical recordings. Its pose estimator is adapted to each participant using that participant's labelled video. Whether a particular day chosen for a later test contributed to training that estimator remains unverified, as does a numerical bound on alignment error. Even when a neural model has never seen a participant or day, the method used to create the movement labels may already have been adapted to that participant. That dependence belongs in the account of what was held out. \href{https://doi.org/10.1038/s41597-022-01280-y}{Peterson et al., 2022}; \href{https://doi.org/10.48324/dandi.000055/0.220127.0436}{AJILE12 release}.

Those sources could contribute different parts of a larger research programme. Their recordings do not share an event, a body or a clock. Combining a macaque's brain--muscle example with a person's brain--pose example would not produce an observed brain--muscle--movement triple. Separate, qualified pretraining could learn useful patterns from each source, followed by a declared transfer whose predictions are checked against actual target measurements. The assumed correspondence would still need testing. The release descriptions identify possible source relationships. They do not establish that the recordings have been acquired or that a model has learned from them. Their exact use conditions, file integrity, units, clocks, earlier exposure and the independence of the chosen partitions still need checking before training.

\section{A world model must account for the body it moves}

Watch a hand close around an object. From outside, the fingers change position, touch begins and the object may move. The body model also needs an inside: muscles and tendons exert forces, joints constrain motion and the load resists it. The proposed world model would relate these events, so that a forecast of movement can be checked against what the declared mechanics and observations support.

Its initial interior consists of specified muscles, tendons, joints and mechanical states. Its exterior consists of supported movement, objects, contacts and environmental conditions. These are deliberately stated parts of a body and a task. Calling the model a world model does not establish complete organs, a body-wide neural recording or access to private sensation. An absent measurement stays absent, however convincing the picture becomes.

The existing component names describe proposed jobs. NeuroTranslate would map declared observations to named targets. MoCap ReComp would join source-labelled movement observations in time and a common physical frame; a collection of estimated video poses does not become that proposed instrument by changing its label. MyoSim Replay would calculate the declared body's mechanical response. Their proposed combination is the Large Movement Model. A neural composite would bring identified recording types and their source-dependent representations into a model. It would not establish that a cell in one person matches a cell in another, or build an observed whole nervous system from unrelated parts.

\subsection{Give every internal value a meaning}

A proposed mechanical state could include joint configuration $q$, generalized velocity $v$, muscle-model activation $a$ and environmental state $c$. Configuration describes how the model is arranged; velocity describes how it is changing. A rotational joint may use radians and radians per second, while a sliding coordinate needs length and length per second. The actual simulator determines which coordinates it uses. Its velocity coordinates need not be identical to ordinary derivatives of every stored position coordinate, so that conversion must be stated.

The model could use permitted earlier observations, task context and a declared control input to predict a range of possible later states. Predicting a range exposes uncertainty that a single animation can hide. During use, the predictor may receive only the inputs allowed for that route. A future measurement used as a training target does not become available before it occurs.

The instruments observe different things. A neural recording supplies its named electrical or imaging signal. EMG supplies electrical activity associated with muscle; a processed EMG amplitude is not automatically a measured activation state or force. Tracking supplies measured or estimated positions, with the camera geometry, calibration and uncertainty attached. A force calculated inside the virtual body remains a calculated force until an independent reference tests it. An availability mask tells the model which values are absent. Filling the numerical array with zero for computation must never be described as having measured zero.

The mechanical accounting can be expressed as
\[
M(q)\dot v+h(q,v)
 =R(q)F_{\mathrm{mt}}+J_c(q)^{\mathsf T}f_c+\tau_{\mathrm{ext}}.
\]
The left side describes the movement changes and the model's declared velocity and gravitational effects. $M$ represents how its mass is distributed and resists acceleration; $\dot v$ describes change in generalized velocity. The term $h$ collects the specified velocity-related and gravitational contributions. On the right, $F_{\mathrm{mt}}$ contains muscle--tendon forces, and $R$ maps their pulls into forces or turning effects on the model's coordinates. Contact forces $f_c$ act through a mapping $J_c^{\mathsf T}$. The final term $\tau_{\mathrm{ext}}$ accounts for other declared external contributions.

Each side must use compatible coordinates and force or turning-effect units for each component. The equation is proposed bookkeeping for a chosen articulated-body simulator. It is not a discovery of all the forces inside a living person. Its geometry, mass distribution, muscle parameters, loads and contact rules need evidence. Established OpenSim methods show how musculoskeletal models can be specified; they do not validate the proposed combination here. \href{https://journals.plos.org/ploscompbiol/article?id=10.1371/journal.pcbi.1006223}{Seth et al., 2018}.

\subsection{Make the visible movement answer to an independent measure}

A predicted path and its reference must describe the same interval, body frame and units. Internal force claims need suitable force measurements as well: different combinations of muscle forces can be compatible with a similar visible path. The tests would examine contacts, allowed joint motion, energy accounting, gravity, friction and sensitivity to numerical solver settings within the claimed domain. A score can improve because the model exploits a numerical correction or has been given future targets. That result would identify a failure in the computation or test, rather than improved forecasting of the intended movement.

The proposed comparison would also remove or replace inputs and constraints one at a time, using matched conditions, to ask what the neural representation, muscle observations, movement prior and physics each contribute. Independent measured outcomes would decide whether those changes help. Later reward-based policy training would need suitable observations of actions, their consequences and rewards, plus a reward calculation checked against those independent outcomes. A dataset of positions alone does not supply all those consequences.

A useful first result could be a forecast for one body region with missing states plainly identified. Predicting it, replaying a recording, controlling a virtual body, intervening in a real body and leaving a person with a lasting skill would still be different achievements. None is established by a more lifelike render.

% R4 full-depth accessible insertion candidate; existing manuscript untouched.
% Inherit the selected master's sans-serif text, code and mathematics fonts.
\subsection{One muscle pulse, two simulated trajectories}

Give the computer hand one muscle command. Does its calculated movement change? That is the question tested here. The installed \texttt{myohand\_r} mechanism represents a human right hand attached to a fixed passive torso. Two copies of the same starting state were compared, with a short excitation added to one of them. No recorded brain activity, muscle electrical signal or measured human movement drove or verified this test. It checks how the simulation responds to a command; the larger biological and training questions remain open. The software was \texttt{myo\_sim} 0.2.3 at commit \texttt{93b0ca8f4ec90c9899ee7f05fee561e9911da91b}, MuJoCo 3.14.0 and NumPy 2.5.3.\cite{myohandMyoSimPinned2026,myohandMuJoCo3140,mujoco2012,myohandNumericalRecord20261105}

\subsubsection{Change the input; keep the starting state}

The compiled hand has 23 hinge coordinates: 23 scalar angles that describe its movable joints. Angles use radians. It also has 39 muscle--tendon actuators, the parts that calculate a muscle-driven pull, and 39 activation states, the internal numbers describing the muscle models' activation. There are zero sensors. A calculated activation is therefore not a measurement of muscle electrical activity, and a calculated force is not a force-instrument reading.

The two branches began with matching complete recorded initial states, including joint angles, velocities and activation states. One branch received no additional excitation. The other received excitation 0.2 at the actuator named \texttt{FDS2\_r}; all other commands stayed zero. Excitation is dimensionless here: 0.2 is a value on the model's control scale, not 0.2 volts, 0.2 newtons or an EMG amplitude. The command drives the muscle model rather than directly requesting a joint angle.

Each branch advanced through 100 steps of $0.002\,\mathrm{s}$, for a nominal $0.200\,\mathrm{s}$ of simulated time. Both used the simulator's native Euler setting, its specified numerical rule for advancing the state. Counting steps from zero, the pulse was present at steps 10 through 59. Those commands apply over the pre-step interval $0.020\leq t<0.120\,\mathrm{s}$. The unpulsed hand could still move under gravity and passive mechanics. Comparing the two trajectories accounts for that shared motion; comparing only the pulsed hand with its starting pose would answer a different question.

A forward evaluation recalculates quantities that depend on the current state, such as the endpoint force. One was made at the beginning and one at the end of each branch. There were therefore 200 explicit integration steps and four explicit forward evaluations. The executed construction path also implies two compilation calls and one internal forward evaluation while the model was being built. These counts do not expose or count every operation inside the native solver.

The built mechanism must be described as it actually was. Leaving its optional \texttt{inertia\_floor} setting at \texttt{None} did not leave all inertias untouched: the native builder changed each diagonal rotational-inertia component of four fixed scaffold bodies from $1$ to $10^{-4}\,\mathrm{kg\,m^2}$. Rotational inertia describes resistance to a change in turning motion. The comparison includes that construction change. No extra physical settings were retuned and no fallback to the alternative \texttt{implicitfast} setting was used. Position, orientation, body scale, contact behaviour and fidelity to a real hand remain uncalibrated by this test.

\subsubsection{Read the difference at the same final state}

Write $P$ for the pulsed branch, $C$ for the control and $T=0.200\,\mathrm{s}$ for the nominal endpoint. The reported comparisons are
\[
\Delta q_{\infty}=\max_{1\leq j\leq23}|q_{P,j}(T)-q_{C,j}(T)|,
\qquad
\Delta a=a_{P,\mathrm{FDS2}}(T)-a_{C,\mathrm{FDS2}}(T),
\]
\[
\Delta F=F_{P,\mathrm{FDS2}}(T)-F_{C,\mathrm{FDS2}}(T).
\]
The first expression subtracts the two final angles for each of the 23 hinges, ignores the sign and reports the largest separation. It does not report the largest separation at any time during the run. The next two expressions keep the sign: pulsed activation minus control activation, and pulsed tendon-actuator force minus control force. Activation has no physical unit; the force uses newtons and the model's sign convention. Both endpoint forces were recalculated at their final states, avoiding a comparison with force values left over from an earlier step.

% Shared by all four proposed R3 inserts; no added packages or font defaults.
\begin{table}[htbp]
\centering
\caption{Final contrasts for one simulated right-hand mechanism. All values compare the pulse branch with its matched control at nominal $T=0.200\,\mathrm{s}$. The angle value is the largest absolute separation across 23 hinges at that endpoint; activation and force retain the pulse-minus-control sign. Activation and force concern the selected \texttt{FDS2\_r} actuator. Force uses the synchronized final forward evaluations and the model's sign convention. These are simulated contrasts, not errors against biological measurements. The displayed digits preserve the recorded numerical values, without establishing biological accuracy.\cite{myohandNumericalRecord20261105}}
\label{tab:myohand-numerical-response}
\begin{tabular}{p{0.53\linewidth}p{0.26\linewidth}p{0.10\linewidth}}
\hline
Quantity & Contrast & Unit \\
\hline
Largest absolute final joint-coordinate separation, $\Delta q_{\infty}$ & $0.5027347840656727$ & rad \\
Selected final activation difference, $\Delta a$ & $0.05205729904534248$ & none \\
Selected final tendon-actuator force difference, $\Delta F$ & $-6.407722136774943$ & N \\
\hline
\end{tabular}
\end{table}


A separate audit checked the saved evidence using standard-library software. It completed 132 checks without a failed check and did not rerun the hand. The required saved numerical channels contained no non-finite values, and the recorded native warning counts were zero. Four saved lists describing the endpoint mass matrices contained positive values. Those lists are numerical diagnostics of how the model resists acceleration. The raw matrices and the operands used for residual calculations were not independently exported, so the audit could check the stored results without independently solving those matrices again.

No contact was observed in this case. At the final control endpoint, the recorded soft-limit excess was $0.0004121829780273552\,\mathrm{rad}$. A soft limit resists motion near a boundary; this observation prevents it being described as an absolute hard stop. The check supplies neither proof of energy conservation nor evidence that changing the numerical settings would give a converged answer. No physiological threshold was used to decide whether a contrast was large enough.

The pulse changed the simulated state under the stated conditions. The table compares two calculated trajectories; it does not measure prediction error against a person. Nor do 200 connected time steps become 200 people or 200 independent experiments. No statistical $p$-value follows from that count. The retained decimal places identify the numerical result without establishing accuracy against biology or repeatability across other cases.

\subsubsection{Keep the unanswered questions visible}

The recorded run used CPU dynamics. Its BLAS numerical-library pool reported one thread, and its dynamics pool size was zero. Temporary parallel work during compilation was not measured. The 263 observed loaded-file identities describe what that bounded inspection saw, rather than every system dependency. Resource use and speed still need analysis. The reported process high-water memory value was 5,187,464 KiB; it can carry history from before an exec transition and cannot be assigned to the hand model alone. It establishes no model-only 3 GiB memory bound. Tool/native timing near 448 ms and timing inside the probe near 491.927 ms use unreconciled boundaries. These values establish neither complete useful end-to-end latency nor a speed improvement, and they do not partition every CPU/GPU cost.

The hand supplies calculated activation, force and motion, with no sensory measurements. Myoception names the programme's proposed sensory side of muscle-driven movement, including touch and body-position feedback. This test does not record those sensory nerve signals or demonstrate a biological feedback loop. It also trains no controller or Large Movement Model and tests no generalization. The wider proposal still has four training-source groups: neural observations, muscle measurements, measured movement and physics-generated examples. Using this generator for training would first require qualified mechanics, controls, calibration, uncertainty, source rights and data lineage, then independent tests against held-out real measurements. More examples and more varied examples remain separate hypotheses. The diagnostic reopens no protected tests or quarantined sources.

The present record has no successful scaling curve or trained body-wide world model. In LINK, eighteen fits and 720 saved epochs still fail both model-admission gates and all five benefit claims. Adding a broader proposal does not undo that outcome. New data must face a new qualified comparison, with the original failures intact and a test strong enough to expose another one.
